\documentclass[11pt,a4paper]{article}
\usepackage[utf8]{inputenc}
\usepackage[T1]{fontenc}
\usepackage[portuguese]{babel}
\usepackage{xcolor}
\usepackage{hyperref}
\usepackage{geometry}
\usepackage{listings}
\usepackage{tcolorbox}
\geometry{margin=2.5cm}
\definecolor{primary}{RGB}{41,128,185}
\definecolor{secondary}{RGB}{44,62,80}
\definecolor{codebg}{RGB}{240,240,240}
\lstset{backgroundcolor=\color{codebg},basicstyle=\ttfamily\small,breaklines=true,frame=single}
\title{\color{secondary}\Huge\textbf{Configuração de Aplicações com ConfigMaps e Secrets}}
\author{\color{primary}DevOps com Kubernetes -- 2026}
\date{}
\begin{document}
\maketitle


\section*{\color{primary}O que você aprenderá nesta página}
\begin{itemize}
\item Separar configuração da imagem da aplicação
\item Usar ConfigMaps para dados não sensíveis
\item Usar Secrets para dados sensíveis e credenciais
\end{itemize}

Uma imagem de contêiner deve ser reutilizável. Se cada ambiente exigir uma imagem diferente apenas para mudar uma URL, uma porta ou uma string de configuração, o processo de publicação fica frágil. O Kubernetes resolve esse problema com recursos declarativos para configuração.

\texttt{ConfigMap} armazena dados de configuração não sensíveis. \texttt{Secret} armazena informações sensíveis, como senhas, tokens e certificados. Ambos podem ser consumidos como variáveis de ambiente ou arquivos montados dentro do contêiner.

\section*{\color{primary}ConfigMaps}

\begin{lstlisting}[language=bash]
apiVersion: v1
kind: ConfigMap
metadata:
  name: app-config
data:
  APP_ENV: production
  API_URL: http://api-svc:3000
\end{lstlisting}

A aplicação pode consumir essas chaves como variáveis de ambiente. Também é possível montar o ConfigMap como volume, útil para arquivos como \texttt{nginx.conf}, \texttt{application.yaml} ou configurações de feature flags.

\begin{lstlisting}[language=bash]
envFrom:
  - configMapRef:
      name: app-config
\end{lstlisting}

\section*{\color{primary}Secrets}

Secrets têm propósito diferente: guardar dados que não devem aparecer no manifesto principal da aplicação. Por padrão, os valores são codificados em Base64, o que não é criptografia. Em clusters reais, é necessário cuidar de RBAC, criptografia em repouso e integração com gerenciadores externos de segredo quando disponível.

\begin{lstlisting}[language=bash]
echo -n 'senha-forte' | base64
\end{lstlisting}

\begin{lstlisting}[language=bash]
apiVersion: v1
kind: Secret
metadata:
  name: db-secret
type: Opaque
data:
  POSTGRES_PASSWORD: c2VuaGEtZm9ydGU=
\end{lstlisting}

Uso no Deployment:

\begin{lstlisting}[language=bash]
env:
  - name: POSTGRES_PASSWORD
    valueFrom:
      secretKeyRef:
        name: db-secret
        key: POSTGRES_PASSWORD
\end{lstlisting}

\section*{\color{primary}Quando usar cada um}

Use ConfigMap para URL de serviço, modo de execução e configurações não sensíveis. Use Secret para senha de banco, token de API, credencial de registry, certificado TLS e chaves privadas. Não coloque senha direto no Deployment e não coloque dado sensível em ConfigMap.

\end{document}
